\begin{explanationbox}

	\textbf{\textcolor{CExplanation}{一句话直觉}}

	线面角是斜线与它在平面上的射影所夹的锐角，其正弦等于垂线段与斜线段的比值。

	\medskip
	\textbf{\textcolor{CExplanation}{核心拆解}}
	\begin{description}[style=nextline, leftmargin=2.8em, labelsep=0.5em, itemsep=0.45em]
		\item[\textbf{要点一（构造垂线）：}] 从斜线上任取一点（非交点），向平面作垂线，垂足与交点连线即得射影。
		\item[\textbf{要点二（确定角）：}] 斜线与射影所成的角就是线面角，在直角三角形中体现。
		\item[\textbf{要点三（正弦公式）：}] $\sin\theta = \dfrac{QH}{PQ}$，其中 $QH$ 是点到面的距离，$PQ$ 是斜线段长度。
	\end{description}

	\medskip
	\textbf{\textcolor{CExplanation}{几何本质（直角三角形）}}

	在直角三角形 $QPH$ 中，$\angle QPH$ 为线面角，其对边 $QH$ 为垂线段，斜边 $PQ$ 为斜线段。因此 $\sin\theta = \dfrac{QH}{PQ}$。

	\medskip
	\textbf{\textcolor{CExplanation}{代数意义（边长比）}}

	线面角的正弦值等于点到平面距离与斜线段长度的比值。这一公式将空间角转化为长度之比，便于计算。

	\medskip
	\textbf{\textcolor{CExplanation}{考点价值（求角度）}}
	\begin{itemize}[leftmargin=*, itemsep=0.5em, topsep=0.4em]
		\item \textbf{考法一（直接计算）：} 给出垂线段和斜线段长度，直接用公式求正弦。
		\item \textbf{考法二（逆向求长）：} 已知正弦和一边，求垂线段或斜线段。
	\end{itemize}

	\medskip
	\textbf{\textcolor{CExplanation}{顿悟点}}

	找到直线的垂足是核心：只要作出垂线，线面角立即转化为直角三角形中的锐角。

	\medskip
	\textbf{\textcolor{CExplanation}{使用场景}}
	\begin{itemize}[leftmargin=*, itemsep=0.5em, topsep=0.4em]
		\item \textbf{场景一（立体几何计算）：} 求直线与平面所成角的正弦。
		\item \textbf{场景二（距离转化）：} 通过线面角正弦可以反求点到平面的距离。
	\end{itemize}
\end{explanationbox}
